\section{Related Work}
\label{sec:related}

\paragraph{Video-language models.}
The recent wave of video-language models pairs a visual encoder with a large language model.
Video-ChatGPT~\cite{maaz2023videochatgpt} averages spatial and temporal features and fine-tunes on video-instruction data.
VideoLLaVA~\cite{lin2023videollava} adds a projection layer that aligns video and image representations before the language model.
LLaVA-Next-Video~\cite{zhang2024llavanextvideo} scales frame resolution and count.
SeViLA~\cite{yu2023sevila} localizes keyframes first, then answers questions from those frames.
Chat-UniVi~\cite{jin2024chatunivi} merges spatial and temporal features at multiple granularities with adaptive tokens.
BLIP~\cite{li2022blip} works at the single-image level, training one model for captioning and VQA. We adopt BLIP as the perception backbone in \ours.
Each of these systems expects the full video at inference time. A growing stream with no known end breaks that expectation. We built \ours for exactly that regime: causal, frame-by-frame input with bounded memory.

\paragraph{Streaming perception.}
Yang \etal~\cite{yang2022real} formalize real-time streaming perception for object detection. Their benchmark penalizes a method whenever it falls behind the input frame rate.
StreamPETR~\cite{wang2023streampetr} carries this idea to 3D detection by propagating temporal features across frames.
Two recent efforts come closer to our problem.
VideoStreaming~\cite{qian2024streaming} proposes a training protocol aware of streaming constraints, and Flash-VStream~\cite{zhang2024flashvstream} adds a memory layer for long video streams.
Both, though, target passive description or offline evaluation. Neither supports real-time interactive Q\&A, and neither reasons about which time window a user's question refers to. \ours fills that gap with explicit temporal scope classification and scope-aware memory filtering.

\paragraph{Video question answering.}
NExT-QA~\cite{xiao2021nextqa} tests causal and temporal reasoning on short clips.
STAR~\cite{wu2021star} probes situated reasoning with compositional questions, and MVBench~\cite{li2024mvbench} covers a broad range of video understanding tasks.
EgoSchema~\cite{mangalam2023egoschema} pushes the time horizon to minutes-long egocentric clips from Ego4D.
MovieChat~\cite{song2024moviechat} tackles long videos by merging frame features into a fixed-size buffer.
One trait unites these benchmarks: the model gets the entire video before it answers. No existing benchmark tests a model that must answer mid-stream, with the future still unseen. LiveQA-Bench, introduced in this paper, fills that gap by requiring causal-only access and temporally-grounded answers on simulated live streams.

\paragraph{Two-stage vision-language pipelines.}
Separating perception from language generation has practical benefits. Each stage uses a frozen pre-trained model, so the pipeline needs no joint training.
BLIP~\cite{li2022blip} and BLIP-2~\cite{li2023blip2} follow this pattern. A vision model produces captions or embeddings in stage one, and a language model reasons over them in stage two.
BLIP-2 pairs a frozen image encoder with Flan-T5~\cite{chung2022flanT5}, an instruction-tuned model covering 1,836 tasks, to answer visual questions without any task-specific training.
Our SFG borrows the same two-stage idea. BLIP extracts per-frame observations, and Flan-T5 fuses them into one coherent answer.

\paragraph{Temporal memory for sequence models.}
How to maintain a fixed-size memory over a potentially infinite input is an old question. Recurrent Memory Transformers~\cite{bulatov2022recurrent} keep a set of memory tokens that persist across segments and carry context forward. Memorizing Transformers~\cite{wu2022memorizing} take a different path: a non-differentiable $k$-nearest-neighbor index over a large external bank.
In the video domain, MovieChat~\cite{song2024moviechat} merges frame features into a buffer, and Ma \etal~\cite{ma2023livelylivechat} compress tokens for long conversations.
\ours departs from these designs in two ways. First, the SKM ranks stored frames by a visual-semantic importance score, not by recency alone. A FIFO buffer drops early events once the window moves forward. Our scoring keeps rare but informative frames. Second, a temporal router sits on top of memory and, at query time, selects only the time slice the question refers to.

\begin{figure*}[!t]
  \centering
  \begin{tikzpicture}[
    node distance=0.6cm and 1.2cm,
    box/.style={draw, rounded corners=4pt, minimum height=0.9cm, minimum width=2.0cm, align=center, font=\small},
    io/.style={draw, rounded corners=4pt, minimum height=0.8cm, minimum width=1.6cm, align=center, font=\small\itshape, dashed},
    arr/.style={-{Stealth[length=5pt]}, thick},
    lbl/.style={font=\scriptsize, midway, fill=white, inner sep=1pt, yshift=3pt},
  ]
    % Row 1: video stream path
    \node[io] (stream) {Live Video\\Stream};
    \node[box, right=1.2cm of stream, fill=blue!8] (enc) {CLIP\\Encoder};
    \node[box, right=1.0cm of enc, fill=blue!8] (skm) {Semantic\\Keyframe\\Memory};
    \node[box, right=1.6cm of skm, fill=green!8] (blip) {BLIP\\Perception};
    \node[box, right=1.4cm of blip, fill=green!8] (flan) {Flan-T5\\Synthesis};
    \node[io, right=1.2cm of flan] (answer) {Real-Time\\Answer};

    % Row 2: query path
    \node[io, below=1.6cm of stream] (query) {User\\Question};
    \node[box, right=1.2cm of query, fill=orange!10] (tqr) {Temporal Query\\Router};

    % Arrows: video stream path
    \draw[arr] (stream) -- (enc) node[lbl, above] {frames};
    \draw[arr] (enc) -- (skm) node[lbl, above] {$v_t$};
    \draw[arr] (skm) -- (blip) node[lbl, above] {keyframes};
    \draw[arr] (blip) -- (flan) node[lbl, above] {captions};
    \draw[arr] (flan) -- (answer) node[lbl, above] {$a$};

    % Arrow: query to TQR
    \draw[arr] (query) -- (tqr) node[lbl, above] {$q$};

    % Arrow: TQR to BLIP via clean L-path (right then up)
    \coordinate (turn) at (blip.south |- tqr.east);
    \draw[arr] (tqr.east) -- (turn) node[lbl, above, pos=0.7] {scope} -- (blip.south);

    % Dashed module groups
    \begin{scope}[on background layer]
      \node[fit=(enc)(skm), draw, dashed, rounded corners=6pt, inner sep=8pt,
            label={[font=\scriptsize]above:SKM}] {};
      \node[fit=(tqr), draw, dashed, rounded corners=6pt, inner sep=8pt,
            label={[font=\scriptsize]above:TQR}] {};
      \node[fit=(blip)(flan), draw, dashed, rounded corners=6pt, inner sep=8pt,
            label={[font=\scriptsize]above:SFG (two-stage)}] {};
    \end{scope}
  \end{tikzpicture}
  \caption{\textbf{Architecture of \ours.}}
  \label{fig:architecture}
\end{figure*}

\begin{figure*}[!t]
  \centering
  \begin{tikzpicture}[
    node distance=0.6cm and 1.8cm,
    box/.style={draw, rounded corners=4pt, minimum height=0.9cm, minimum width=2.2cm, align=center, font=\small},
    arr/.style={-{Stealth[length=5pt]}, thick},
    lbl/.style={font=\scriptsize, midway, fill=white, inner sep=1pt, yshift=3pt},
    iolbl/.style={font=\small\itshape},
  ]
    \node[iolbl] (qin) {$q$};
    \node[box, right=1.6cm of qin, fill=orange!10] (qenc) {Keyword\\Matcher};
    \node[box, right=2.0cm of qenc, fill=orange!10] (scope) {Scope\\Selector};
    \node[box, right=2.0cm of scope, fill=orange!10] (filter) {Memory\\Filter};
    \node[iolbl, right=1.6cm of filter] (fout) {$\mathcal{F}$};

    \draw[arr] (qin) -- (qenc);
    \draw[arr] (qenc) -- (scope) node[lbl, above] {per-scope scores};
    \draw[arr] (scope) -- (filter) node[lbl, above] {scope $c$};
    \draw[arr] (filter) -- (fout);
  \end{tikzpicture}
  \caption{\textbf{TQR pipeline.}}
  \label{fig:tqr}
\end{figure*}
